\chapter*{\centering ABSTRACT}
\label{ch:abstract}
\addcontentsline{toc}{chapter}{Abstract}

The Kubernetes Horizontal Pod Autoscaler (HPA) resizes a deployment only
after a chosen utilization metric crosses a fixed threshold, and it cannot
use experience from past traffic to anticipate future load. This report
describes a capstone project that trained a Proximal Policy Optimization
(PPO) agent directly against a live, single-node Minikube cluster, with no
simulator standing in for the cluster at any stage, and compared it against
two HPA configurations on traffic none of the controllers had seen before.
The agent observes eleven normalized signals covering CPU, memory, request
rate, and latency, together with their short-horizon trends; its raw action
passes through a deterministic safety layer -- a replica-bound clamp, a
pending-scale gate, and an asymmetric scale-down cooldown -- before it ever
reaches the Kubernetes API, and an uncertain policy is resolved by a
safety-aware tie-break rule that can never let scale-down win a tie. The
agent trained for 12,288 live decision steps across 48 sessions, during
which explained variance of the value function rose from 0.0125 at the
first update to a mean of 0.703 over the last five updates (peaking at
0.799), computed directly from the training campaign's own update log. Two
of the 48 planned sessions stopped early when every ready replica of the
deployment briefly dropped to zero under correlated pod unhealthiness; both
resumed from their last checkpoint and completed successfully, and we
report this as a genuine property of running training against a real
cluster rather than omitting it. We then ran a pre-registered, 27-run
held-out evaluation of this policy against an HPA configuration using CPU
and memory and a CPU-only HPA baseline (the same 260\% CPU target and 2-8
replica bounds as every other arm, not Kubernetes' unmodified factory
settings), across
three held-out traffic profiles, followed by a separately labelled,
exploratory nine-run stress test at training-level peak loads. All 27
pre-registered runs were valid; no controller was declared an overall
winner, because our decision rule required a controller to win on a metric
in all three repeats of a profile, and cost rankings differed by profile:
the CPU-only HPA was cheapest on two of three profiles, while the PPO
policy was cheaper than the CPU-and-memory HPA on two profiles but more
expensive on the third. The exploratory stress test likewise found no
controller that was better across every profile. We report these results,
together with the operational incidents and limitations encountered along
the way, as evidence about what a live-trained RL autoscaler can and cannot
be shown to do on a small cluster within one capstone's timeline.

\vspace{0.8cm}
\noindent \textbf{Keywords:} Kubernetes, autoscaling, reinforcement learning,
Proximal Policy Optimization, Horizontal Pod Autoscaler, cloud resource
management, Minikube, pre-registration
